% Profielbestuur
\section{Profielbestuur}
\subsection{Gebruikersgeval}

\begin{figure}[htbp]
    \centering
    \includegraphics[width=\textwidth]{figures/usecase/profile.png}
    \caption{Gebruikersgeval vir Profielbestuur}
    \label{fig:profielbestuur_use_case}
\end{figure}

Volgens FREQ-USR-201, FREQ-USR-202 en FREQ-USR-203 moet gebruikers hul profiel kan besigtig (\autoref{fig:profile_screen}) en hul koshuis en kampus ligging, allergene en wagwoord verander deur 'n profielwysigingsproses (\autoref{fig:profile_edit_screen}). Verder moet die gebruiker hul huidige rekeningbalans en transaksie geskiedenis kan sien volgens FREQ-MOB-203 en FREQ-MOB-204.

\subsection{Scenario}

\textbf{Akteurs:} Klient\\
\textbf{Voorvereistes:} Klient is ingeteken\\
\textbf{Prekondisie:} Klientprofiel is geskep\\
\textbf{Postkondisie:} Klientprofiel word gewysig\\
\textbf{Inset:} Inligting wat gewysig moet word\\
\textbf{Uitset:} Sukses- of foutboodskap\\
\textbf{Skerms van Toepassing:} 
\begin{itemize}
    \item \autoref{fig:profile_screen} (Profielbesigtig skerm)
    \item \autoref{fig:profile_edit_screen} (Profielwysiging skerm)
\end{itemize}

\textbf{Hoofvloei:}
\begin{enumerate}
    \item Gebruiker navigeer na profiel bladsy (\autoref{fig:profile_screen})
    \item Stelsel toon huidige profiel inligting (\autoref{fig:profile_screen})
    \item Stelsel toon huidige rekeningbalans
    \item Stelsel toon transaksiegeskiedenis (regs bo)
    \item Gebruiker kies om profiel te wysig (\autoref{fig:profile_edit_screen})
    \item Stelsel maak velde bewerkbaar
    \item Gebruiker wysig gewenste inligting (koshuis, kampus, allergieë, wagwoord)
    \item Gebruiker klik "Stoor Veranderinge"
    \item Stelsel valideer gewysigde inligting
    \item Stelsel werk profiel by
    \item Stelsel toon sukses boodskap
    \item Stelsel toon bygewerkte profiel inligting (\autoref{fig:profile_screen})
    \item Gebruikersgeval eindig
\end{enumerate}

\textbf{Alternatiewe vloei A1: Gebruiker kyk net profiel}
\begin{enumerate}
    \item[5a.] Gebruiker kies nie om te wysig nie
    \item Gebruikersgeval eindig
\end{enumerate}

\textbf{Alternatiewe vloei A2: Wagwoord wysiging}
\begin{enumerate}
    \item[7a.] Gebruiker kies om wagwoord te wysig
    \item Stelsel vra huidige wagwoord
    \item Stelsel vra nuwe wagwoord (twee keer)
    \item Stelsel valideer huidige wagwoord
    \item Gaan verder met stap 8
\end{enumerate}

\textbf{Uitsonderingsvloei E1: Ongeldige of uitgelate wysigings}
\begin{enumerate}
    \item[9a.] Stelsel identifiseer ongeldige of uitgelate inligting
    \item Stelsel toon foutboodskappe
    \item Gaan terug na stap 7
\end{enumerate}

